% This document was generated by an LLM.

\documentclass{essay}

\title{An Obvious Theorem, Proved Twice\\[2pt]
\large On the convergence of $r^{n}$ to $0$ for $|r| < 1$}
\author{}
\date{}

\begin{document}

\maketitle

\begin{abstract}
\noindent
The assertion that $r^{n} \to 0$ whenever $|r| < 1$ is among the first facts a
student meets about sequences, and it is usually treated as self-evident: repeated
multiplication by something small ought to produce something small. This note takes
the statement seriously. We isolate exactly which axioms of the real number system
the result depends on, prove it twice by genuinely different mechanisms, and exhibit
an ordered field in which the statement is \emph{false}. The purpose is not the
theorem, which is not in doubt, but the discipline: to display what it looks like to
discharge every obligation in an argument whose conclusion nobody disputes.
\end{abstract}

\tableofcontents

\section{Statement and standing conventions}

\begin{theorem}\label{thm:main}
Let $r \in \mathbb{C}$ with $|r| < 1$. Then the sequence $(r^{n})_{n \in \mathbb{N}}$
converges to $0$; that is,
\[
  \forall \varepsilon > 0 \;\; \exists N \in \mathbb{N} \;\; \forall n \geq N :
  \quad |r^{n} - 0| < \varepsilon .
\]
\end{theorem}

We write $\mathbb{N} = \{0, 1, 2, \dots\}$. Throughout, $\mathbb{R}$ denotes a complete
ordered field, and $|\cdot|$ denotes the usual absolute value on $\mathbb{R}$ or the
modulus on $\mathbb{C}$, as context requires. The only properties of $\mathbb{R}$ that
will be invoked are the ordered-field axioms together with the following.

\begin{definition}[Least upper bound axiom]\label{def:lub}
Every nonempty subset of $\mathbb{R}$ that is bounded above possesses a least upper
bound (supremum) in $\mathbb{R}$.
\end{definition}

Every appeal to ``completeness'' below is an appeal to Definition~\ref{def:lub},
directly or through one of its consequences. This bookkeeping is the point of the
essay: the theorem is not a fact about arithmetic, it is a fact about the
completeness of $\mathbb{R}$, and Section~\ref{sec:counterexample} makes that
concrete.

\section{Reduction to the real, positive case}

\begin{lemma}[Multiplicativity of the modulus on powers]\label{lem:mult}
For all $r \in \mathbb{C}$ and all $n \in \mathbb{N}$ we have $|r^{n}| = |r|^{n}$.
\end{lemma}

\begin{proof}
Induction on $n$. For $n = 0$, both sides equal $1$, since $r^{0} = 1$ and
$|1| = 1$. Suppose $|r^{n}| = |r|^{n}$ for some $n \in \mathbb{N}$. Then, using
multiplicativity of the modulus,
\[
  |r^{n+1}| = |r^{n} \cdot r| = |r^{n}| \cdot |r| = |r|^{n} \cdot |r| = |r|^{n+1}.
\]
By induction the identity holds for all $n \in \mathbb{N}$.
\end{proof}

\begin{proposition}[Reduction]\label{prop:reduction}
Theorem~\ref{thm:main} follows from the special case in which $r$ is real and
$0 < r < 1$.
\end{proposition}

\begin{proof}
Set $p := |r|$, so $p \in [0,1)$ by hypothesis. By Lemma~\ref{lem:mult},
\[
  |r^{n} - 0| = |r^{n}| = p^{n} = |p^{n} - 0|
  \qquad \text{for every } n \in \mathbb{N},
\]
so the statement ``$r^{n} \to 0$'' and the statement ``$p^{n} \to 0$'' unfold to
the identical quantified assertion, and each holds if and only if the other does.

If $p = 0$, then $r = 0$ and $r^{n} = 0$ for all $n \geq 1$; the definition of
convergence is satisfied with $N = 1$ for every $\varepsilon > 0$, since
$|r^{n} - 0| = 0 < \varepsilon$. Otherwise $p \in (0,1)$, which is the special case
assumed available.
\end{proof}

For the remainder of the essay we fix $p \in \mathbb{R}$ with $0 < p < 1$ and prove
that $p^{n} \to 0$.

\begin{lemma}[Positivity of the powers]\label{lem:positive}
$p^{n} > 0$ for every $n \in \mathbb{N}$.
\end{lemma}

\begin{proof}
Induction. $p^{0} = 1 > 0$. If $p^{n} > 0$, then $p^{n+1} = p^{n} \cdot p$ is a
product of two strictly positive elements of an ordered field, hence strictly
positive.
\end{proof}

\section{First proof: monotone convergence and the fixed-point equation}

The strategy is to separate two questions that are easy to conflate: \emph{whether}
the limit exists, and \emph{what} its value is. Completeness answers the first;
an algebraic identity answers the second.

\subsection{The tools}

\begin{lemma}[Monotone convergence theorem, decreasing form]\label{lem:mct}
Let $(a_{n})_{n \in \mathbb{N}}$ be a sequence of real numbers that is nonincreasing
$(a_{n+1} \leq a_{n}$ for all $n)$ and bounded below. Then $(a_{n})$ converges, and
\[
  \lim_{n \to \infty} a_{n} = \inf \{ a_{n} : n \in \mathbb{N} \}.
\]
\end{lemma}

\begin{proof}
Let $S = \{a_{n} : n \in \mathbb{N}\}$, a nonempty set bounded below. Then $-S$ is
nonempty and bounded above, so by Definition~\ref{def:lub} it has a supremum;
setting $L := -\sup(-S)$ gives the infimum of $S$. We claim $a_{n} \to L$.

Let $\varepsilon > 0$. Since $L + \varepsilon > L$ and $L$ is the \emph{greatest}
lower bound of $S$, the number $L + \varepsilon$ is not a lower bound of $S$; hence
there exists $N \in \mathbb{N}$ with $a_{N} < L + \varepsilon$. For any $n \geq N$,
monotonicity gives $a_{n} \leq a_{N}$ (a trivial induction on $n - N$), while $L$
being a lower bound gives $L \leq a_{n}$. Therefore
\[
  L - \varepsilon < L \leq a_{n} \leq a_{N} < L + \varepsilon,
\]
so $|a_{n} - L| < \varepsilon$. As $\varepsilon > 0$ was arbitrary, $a_{n} \to L$.
\end{proof}

\begin{lemma}[Shift invariance of limits]\label{lem:shift}
If $a_{n} \to L$, then $a_{n+1} \to L$.
\end{lemma}

\begin{proof}
Let $\varepsilon > 0$ and choose $N$ with $|a_{n} - L| < \varepsilon$ for all
$n \geq N$. If $n \geq N$ then $n + 1 \geq N$, so $|a_{n+1} - L| < \varepsilon$.
\end{proof}

\begin{lemma}[Scalar multiples]\label{lem:scalar}
If $a_{n} \to L$ and $c \in \mathbb{R}$, then $c\,a_{n} \to cL$.
\end{lemma}

\begin{proof}
If $c = 0$ the claim is immediate. Otherwise let $\varepsilon > 0$ and choose $N$
with $|a_{n} - L| < \varepsilon / |c|$ for all $n \geq N$; then
$|c a_{n} - cL| = |c|\,|a_{n} - L| < \varepsilon$ for such $n$.
\end{proof}

\begin{lemma}[Uniqueness of limits]\label{lem:unique}
A sequence of real numbers has at most one limit.
\end{lemma}

\begin{proof}
Suppose $a_{n} \to L$ and $a_{n} \to M$ with $L \neq M$. Put
$\varepsilon := |L - M| / 2 > 0$ and choose $N_{1}, N_{2}$ witnessing convergence to
$L$ and to $M$ respectively for this $\varepsilon$. For $n \geq \max(N_{1}, N_{2})$
the triangle inequality gives
\[
  |L - M| \leq |L - a_{n}| + |a_{n} - M| < \varepsilon + \varepsilon = |L - M|,
\]
a contradiction.
\end{proof}

\subsection{The proof}

\begin{proof}[First proof of Theorem~\ref{thm:main}]
By Proposition~\ref{prop:reduction} it suffices to show $p^{n} \to 0$ for
$0 < p < 1$.

\medskip
\noindent\textbf{Step 1: the sequence is strictly decreasing and bounded below.}
By Lemma~\ref{lem:positive}, $p^{n} > 0$ for every $n$, so $0$ is a lower bound.
Moreover, multiplying the strict inequality $p < 1$ by the strictly positive number
$p^{n}$ preserves it, giving
\[
  p^{n+1} = p^{n} \cdot p < p^{n} \cdot 1 = p^{n}.
\]

\medskip
\noindent\textbf{Step 2: the sequence converges.}
By Step 1 and Lemma~\ref{lem:mct}, there exists $L \in \mathbb{R}$ with
$p^{n} \to L$. Since $0$ is a lower bound for the terms and $L$ is their infimum,
$L \geq 0$.

\medskip
\noindent\textbf{Step 3: the limit satisfies $L = pL$.}
Apply Lemma~\ref{lem:shift} to the convergent sequence $(p^{n})$: the shifted
sequence $(p^{n+1})$ converges to $L$ as well. Apply Lemma~\ref{lem:scalar} with
$c = p$ to the same sequence: since $p^{n+1} = p \cdot p^{n}$, the sequence
$(p^{n+1})$ converges to $pL$. By Lemma~\ref{lem:unique} a sequence cannot have two
distinct limits, so
\[
  L = pL.
\]

\medskip
\noindent\textbf{Step 4: conclude.}
Rearranging, $L(1 - p) = 0$. In a field there are no zero divisors, and
$1 - p \neq 0$ because $p < 1$. Hence $L = 0$, i.e. $p^{n} \to 0$.
\end{proof}

\begin{remark}[The order of Steps 2 and 3 is not negotiable]\label{rem:error}
The identity $L = pL$ is meaningless until $L$ has been shown to exist. This is the
standard error in this proof, and it is not a pedantic one: the manipulation
``$L = \lim q^{n+1} = q \lim q^{n} = qL$, hence $L(1-q) = 0$, hence $L = 0$'' applies
verbatim to $q = 2$ and yields the false conclusion that $2^{n} \to 0$. What fails
for $q = 2$ is precisely Step 2, and nothing else. An argument whose validity turns
entirely on a step one is tempted to omit is an argument worth writing out.
\end{remark}

\begin{remark}[What Step 3 really is]
Step 3 says the limit must be a fixed point of the map $x \mapsto px$. That map has
$0$ as its unique fixed point exactly because $p \neq 1$. The first proof is thus a
rigidity argument: completeness produces a limit, and the algebra of the recursion
leaves it no room to be anything but $0$.
\end{remark}

\section{Second proof: Bernoulli's inequality and the Archimedean property}

The first proof establishes convergence without ever producing an $N$ explicitly.
The second proof does the opposite: it constructs an explicit bound on $p^{n}$ and
reads $N$ off from it. It uses completeness only through the Archimedean property.

\subsection{The tools}

\begin{lemma}[Bernoulli's inequality]\label{lem:bernoulli}
Let $h \in \mathbb{R}$ with $h \geq -1$. Then $(1+h)^{n} \geq 1 + nh$ for every
$n \in \mathbb{N}$.
\end{lemma}

\begin{proof}
Induction on $n$. For $n = 0$ both sides equal $1$. Assume $(1+h)^{n} \geq 1 + nh$.
Since $1 + h \geq 0$, multiplying this inequality by $1 + h$ preserves its direction:
\[
  (1+h)^{n+1} = (1+h)^{n}(1+h) \geq (1 + nh)(1 + h)
  = 1 + (n+1)h + nh^{2} \geq 1 + (n+1)h,
\]
the last step because $nh^{2} \geq 0$. This completes the induction.
\end{proof}

\begin{lemma}[Archimedean property]\label{lem:archimedes}
For every $x \in \mathbb{R}$ there exists $N \in \mathbb{N}$ with $N > x$.
\end{lemma}

\begin{proof}
Suppose not. Then $\mathbb{N} \subseteq \mathbb{R}$ is nonempty and bounded above by
$x$, so by Definition~\ref{def:lub} it has a supremum $s$. Since $s - 1 < s$ and $s$
is the least upper bound, $s - 1$ is not an upper bound: there is $m \in \mathbb{N}$
with $m > s - 1$. But then $m + 1 \in \mathbb{N}$ and $m + 1 > s$, contradicting
that $s$ is an upper bound for $\mathbb{N}$.
\end{proof}

\subsection{The proof}

\begin{proof}[Second proof of Theorem~\ref{thm:main}]
Again it suffices to treat $p$ with $0 < p < 1$.

\medskip
\noindent\textbf{Step 1: a change of variable.}
Define
\[
  h := \frac{1}{p} - 1 = \frac{1 - p}{p}.
\]
Since $0 < p < 1$, both numerator and denominator are strictly positive, so $h > 0$;
and by construction $1 + h = 1/p$, that is,
\[
  p = \frac{1}{1 + h}, \qquad h > 0.
\]

\medskip
\noindent\textbf{Step 2: an explicit upper bound for $p^{n}$.}
Raising the previous identity to the $n$-th power and applying
Lemma~\ref{lem:bernoulli},
\[
  p^{n} = \frac{1}{(1+h)^{n}} \leq \frac{1}{1 + nh}
  \qquad \text{for all } n \in \mathbb{N},
\]
where the inequality is obtained from $(1+h)^{n} \geq 1 + nh > 0$ by taking
reciprocals, which reverses inequalities between strictly positive quantities.
Combining with Lemma~\ref{lem:positive},
\begin{equation}\label{eq:bound}
  0 < p^{n} \leq \frac{1}{1 + nh} < \frac{1}{nh}
  \qquad \text{for all } n \geq 1 .
\end{equation}

\medskip
\noindent\textbf{Step 3: choose $N$.}
Let $\varepsilon > 0$ be given. By Lemma~\ref{lem:archimedes} applied to the real
number $1 / (h\varepsilon)$, there exists $N \in \mathbb{N}$ with
\[
  N > \frac{1}{h \varepsilon} .
\]
Note $N \geq 1$, since $1/(h\varepsilon) > 0$. For any $n \geq N$ we have
$nh \geq Nh > 1/\varepsilon > 0$, hence $1/(nh) < \varepsilon$, and so
by \eqref{eq:bound},
\[
  |p^{n} - 0| = p^{n} < \frac{1}{nh} < \varepsilon .
\]
Since $\varepsilon > 0$ was arbitrary, $p^{n} \to 0$.
\end{proof}

\begin{remark}[The two proofs are genuinely different]
The first proof is non-constructive in an essential way: it produces the limit as a
supremum supplied by Definition~\ref{def:lub} and never exhibits, for a given
$\varepsilon$, any particular index beyond which the terms are small. The second
proof is quantitative. It yields the explicit and reusable estimate
\[
  p^{n} \leq \frac{1}{1 + n\left(\tfrac{1}{p} - 1\right)} = O(n^{-1}),
\]
which is a crude bound---the truth is geometric decay---but a bound one can
sharpen. Replacing Bernoulli's inequality by the stronger binomial estimate
$(1+h)^{n} \geq \binom{n}{2} h^{2}$, valid for $h > 0$ and $n \geq 2$, gives
$p^{n} = O(n^{-2})$ and hence the strictly stronger conclusion $n p^{n} \to 0$,
which does not follow from the first proof without further work.
\end{remark}

\section{Why the theorem is not obvious: a countermodel}
\label{sec:counterexample}

An argument's necessity is best demonstrated by removing a hypothesis and watching
the conclusion fail. Both proofs above used completeness: the first through the
monotone convergence theorem, the second through the Archimedean property. Neither
is a consequence of the ordered-field axioms alone, and without them the theorem is
false.

Consider the field $\mathbb{R}(t)$ of rational functions with real coefficients,
ordered by declaring $f > 0$ if and only if $f(x) > 0$ for all sufficiently large
real $x$. One verifies without difficulty that this is a total order compatible with
the field operations, so $\mathbb{R}(t)$ is an ordered field containing $\mathbb{R}$
as an ordered subfield. In it, the element $t$ satisfies $t > n$ for every
$n \in \mathbb{N}$; consequently $\varepsilon_{0} := 1/t$ satisfies
\[
  0 < \varepsilon_{0} < \frac{1}{n} \qquad \text{for every } n \geq 1 .
\]
This field is therefore non-Archimedean: Lemma~\ref{lem:archimedes} fails, and with
it Definition~\ref{def:lub}.

Now take $p = 1/2$, an element of $\mathbb{R}(t)$ satisfying $0 < p < 1$. For every
$n \in \mathbb{N}$ the quantity $2^{-n}$ is an ordinary positive real number, while
$\varepsilon_{0}$ is strictly smaller than every positive real; hence
\[
  p^{n} = 2^{-n} > \varepsilon_{0} \qquad \text{for every } n \in \mathbb{N}.
\]
Thus the decreasing sequence $(p^{n})$ is bounded below by the strictly positive
element $\varepsilon_{0}$, and no matter how the notion of limit is formulated in
the order topology of $\mathbb{R}(t)$, the sequence does not converge to $0$: the
neighbourhood $(-\varepsilon_{0}, \varepsilon_{0})$ of $0$ contains no term of the
sequence at all. In fact $(p^{n})$ has no limit whatsoever, since the set of its
lower bounds---which includes every positive infinitesimal---has no greatest
element.

The moral is exact. The arithmetic of $p^{n}$ is identical in $\mathbb{R}$ and in
$\mathbb{R}(t)$: the same recursion, the same monotonicity, the same positivity.
Everything that ``obviously'' makes the powers shrink to nothing is present in the
countermodel. What is absent is completeness. The theorem is therefore not a
statement about multiplication making things smaller; it is a statement about
$\mathbb{R}$ having no gaps for the sequence to hide in.

\section{Summary of dependencies}

\begin{center}
\begin{tabular}{ll}
\hline
\textbf{First proof} & \textbf{Second proof} \\
\hline
Ordered-field axioms & Ordered-field axioms \\
Least upper bound axiom & Least upper bound axiom \\
\quad via monotone convergence (Lem.~\ref{lem:mct})
  & \quad via the Archimedean property (Lem.~\ref{lem:archimedes}) \\
Uniqueness of limits (Lem.~\ref{lem:unique}) & Bernoulli's inequality
  (Lem.~\ref{lem:bernoulli}) \\
Algebra of limits (Lem.~\ref{lem:shift},~\ref{lem:scalar}) & Induction \\
Induction & \\
\hline
Non-constructive & Constructive: $N > \dfrac{p}{(1-p)\varepsilon}$ \\
\hline
\end{tabular}
\end{center}

\medskip

\noindent
Neither proof appeals to logarithms, exponentials, power series, the ratio test, or
any part of the theory of geometric series. This is deliberate. All of those tools
are downstream of Theorem~\ref{thm:main}: the convergence of $\sum r^{n}$ for
$|r| < 1$, the radius-of-convergence machinery, and the analytic construction of
$\log$ each rest on the fact proved here. A proof of Theorem~\ref{thm:main} that
invokes them is not a proof but a circle, however impressive its vocabulary. The two
arguments above are, in a precise sense, as deep as this statement is allowed to go.

\end{document}
